% Generated by q2_2_accuracy_audit.py; the 11.046 s schedule remains the main Q2(2)/Q4 baseline.
\textbf{Accuracy audit and guarded variant.} Using the same $1^\circ$ independent
bearing-error model as Q2(1), the time-first schedule has RMSE values of
4.640, 2.695, and 1.995 m for T1--T3.
The largest error occurs at T1 because its crossing angle is only
17.65$^\circ$. We therefore retain the
$11.046$ s schedule as the time-objective solution, but additionally report a
$\gamma\ge30^\circ$ guarded variant. Its max start increases by only
0.603 s to $11.649$ s, while T1's crossing angle rises to
108.92$^\circ$ and its RMSE falls to
1.708 m; the worst-alert RMSE falls from
4.640 to 2.873 m.
T2 changes from 2.695 to 2.873 m, so this is a
speed--accuracy trade-off, not a uniform improvement. The remaining metre-level errors
mainly reflect the 40.8--90.8 m station--target distances,
unlike the close-range $90^\circ$ geometry specifically optimized in Q2(1)
($0.14362$ m).

\begin{table}[!htb]\centering\footnotesize\setlength{\tabcolsep}{3.5pt}
\caption{Q2(2) accuracy audit. Each alert cell reports crossing angle $\gamma$ (deg) /
RMSE (m); the time-first schedule remains the main time-objective solution.}
\label{tab:q22accuracy}
\begin{tabular}{lccccc}
\toprule
Schedule & max start (s) & T1 $\gamma$/RMSE & T2 $\gamma$/RMSE & T3 $\gamma$/RMSE & worst RMSE (m) \\
\midrule
Time-first ($\gamma\ge1^\circ$) & 11.046 &
17.65/4.640 &
50.23/2.695 &
59.21/1.995 & 4.640 \\
Accuracy-guarded ($\gamma\ge30^\circ$) & 11.649 &
108.92/1.708 &
31.80/2.873 &
59.21/1.995 & 2.873 \\
\bottomrule
\end{tabular}
\end{table}
